\documentclass[a4paper]{article}
\usepackage{amsfonts}
\usepackage{amsmath}
\usepackage{amssymb}
\usepackage{graphicx}
\DeclareMathOperator{\Bgtan}{Bgtan}

\begin{document}

\noindent \large Auteur: Karanjot Singh \\
\noindent \large Studierichting: Informatica \\
\noindent \large Oefeningenreeks 6A nr. 6.15 \\

\normalsize

$\displaystyle\int_{0}^{1} \dfrac{dx}{x^2 + 4}$ \\

We gebruiken de substitutie $x = 2 \tan \theta$, dus $dx = 2 \sec^2 \theta\, d\theta$ \\

Pas de grenzen aan: \\

Als $x = 0$, dan $\theta = 0$ \\

Als $x = 1$, dan $1 = 2 \tan \theta \Rightarrow \theta = \operatorname{Bgtan}\left( \dfrac{1}{2} \right)$ \\

Vervang in de integrand: \\

$\displaystyle\int_{0}^{1} \dfrac{dx}{x^2 + 4} = \displaystyle\int_{0}^{\operatorname{Bgtan}\left(\tfrac{1}{2}\right)} \dfrac{2 \sec^2 \theta\, d\theta}{4 \tan^2 \theta + 4}$ \\

$= \displaystyle\int_{0}^{\operatorname{Bgtan}\left(\tfrac{1}{2}\right)} \dfrac{2 \sec^2 \theta\, d\theta}{4(\tan^2 \theta + 1)}$ \\

$= \displaystyle\int_{0}^{\operatorname{Bgtan}\left(\tfrac{1}{2}\right)} \dfrac{2 \sec^2 \theta\, d\theta}{4 \sec^2 \theta}$ \\

$= \displaystyle\int_{0}^{\operatorname{Bgtan}\left(\tfrac{1}{2}\right)} \dfrac{2}{4}\, d\theta = \dfrac{1}{2} \displaystyle\int_{0}^{\operatorname{Bgtan}\left(\tfrac{1}{2}\right)} d\theta$ \\

$= \dfrac{1}{2} \left( \operatorname{Bgtan}\left( \dfrac{1}{2} \right) - 0 \right)$ \\

$= \dfrac{1}{2} \operatorname{Bgtan}\left( \dfrac{1}{2} \right)$ \\

Dus, $\displaystyle\int_{0}^{1} \dfrac{dx}{x^2 + 4} = \dfrac{1}{2} \operatorname{Bgtan}\left( \dfrac{1}{2} \right)$ \\

\begin{thebibliography}{999}
	%\bibitem{a} Referentie
\end{thebibliography}

\end{document}
